\documentclass[12pt]{article}

\newif\ifsol
\soltrue

\include{macros}

\begin{document}
\begin{center}
\Large{\textbf{CAS CS 538.  \ifsol Solutions to \fi Problem Set 7}}\\
\smallskip
\large{\textbf{Due electronically via Gradescope, Monday March 30, 2026 11:59pm
}}\\
\large{\textbf{Om Khade, U51801771}}
\end{center}

\noindent
 

\begin{problem} (25 points)
Consider the following MAC (a variant of this was used for Wi\nobreakdash-Fi encryption in 802.11b WEP).
Let $F$ be a PRF defined over $(\K, \R, \X)$ where $\X := \{0, 1\}^{32}$.
Let CRC32 be any efficiently computable function mapping $m \in \{0, 1\}^{\le \ell}$ to a 32-bit string (in 802.11b, it was a simple and popular error-detecting code meant to detect random errors).
Define the following MAC system $(S, V)$:
\begin{align*}
	S(k, m) &:= \left\{r \samples \R, t \gets F(k, r) \oplus \text{CRC32}(m), \text{output }(r, t)\right\}, \\
	V(k, m, (r, t)) &:= \left\{\textsf{accept}\text{ if }t = F(k, r) \oplus \text{CRC32}(m)\text{ and }\textsf{reject}\text{ otherwise}\right\}.
\end{align*}
Show that this MAC system is insecure.
\end{problem} 
  
\begin{solution}
The main issue with this MAC setup is the separation of the nonce $r$ from the msg $m$ via XOR. Basically, the MAC reveals enough info about the hidden value $t = F(k, r)$ to allow an adv. to compute a new valid MAC for a different msg using the same nonce.

\subsection*{Adversary $\A$}
Let $\A$ be an efficent adv attacking the defined MAC protocol:
\begin{itemize}
    \item $\A$ chooses some rnd msg $m_1 \in \M$.
    \item $\A$ queries the signing oracle $S(k, m_1)$, and gets $(r_1, t_1)$.
    \item $\A$ then chooses a rnd msg $m_2 \in \M$ such that $m_2 \neq m_1$.
    \item $\A$ then computes $t_2$:
    \begin{equation*}
        t_2 = t_1 \oplus \text{CRC32}(m_1) \oplus \text{CRC32}(m_2)
    \end{equation*}
    \item $\A$ outputs the candidate forgery $(m_2, (r_1, t_2))$.
\end{itemize}

\subsection*{Analysis}
Note that, in order for $\A$ to win, it must achieve the following properties:
\begin{enumerate}
    \item The msg, $m_2$, was never queried to the signing oracle.
    \item The verification algorithm accepts the pair: $V(k, m_2, (r_1, t_2)) = \textsf{ACCEPT}$.
\end{enumerate}

\paragraph{Freshness} By construction, $\A$ only queried $m_1$. Since $m_2 \neq m_1$, freshness is guaranteed.

\paragraph{Validity} From the signing query, $\A$ knows that $t_1 = F(k, r_1) \oplus \text{CRC32}(m_1)$. 
It is now possible to rearrange this, such that we can isolate the PRF output:
\begin{equation*}
    F(k, r_1) = t_1 \oplus \text{CRC32}(m_1)
\end{equation*}
Now, consider the verification equation for the forgery $(m_2, (r_1, t_2))$. The verifier checks if:
\begin{equation*}
    t_2 \stackrel{?}{=} F(k, r_1) \oplus \text{CRC32}(m_2)
\end{equation*}
By subbing the value of $t_2$, computed by $\A$:
\begin{equation*}
    \begin{aligned}
        t_2 &= (t_1 \oplus \text{CRC32}(m_1)) \oplus \text{CRC32}(m_2) \\
            &= F(k, r_1) \oplus \text{CRC32}(m_2)
    \end{aligned}
\end{equation*}
The equation holds perfectly. Therefore, $V$ will always output $\text{ACCEPT}$.

\subsection*{End}
Since $\A$ wins the game with probability 1:
\begin{equation*}
    \MACadv[\A, \text{MAC}] = 1
\end{equation*}
Since 1 is not negligible, this proves that the MAC protocol is insecure. $\qed$
\end{solution}


  
\begin{problem} (30 points)
We want to build a MAC system $\I$ using two MAC system $\I_1 = (S_1, V_1)$ and $\I_2 = (S_2, V_2)$, so that if at some time one of $\I_1$ or $\I_2$ is broken (but not both) then $\I$ is still secure.
Put another way, we want to construct $\I$ from $\I_1$ and $\I_2$ such that $\I$ is secure if either $\I_1$ or $\I_2$ is secure.
This is also called a \emph{MAC combiner}.

For this problem, suppose that $\I_1$ and $\I_2$ are deterministic MAC systems (see the definition on textbook page 217 in Section 6.1), and that both have tag space $\{0, 1\}^n$.
Define the deterministic MAC system $\I = (S, V)$, where
\[
	S((k_1, k_2), m) := S_1(k_1, m) \oplus S_2(k_2, m).
\]
Show that $\I$ is a MAC combiner of $\I_1$ and $\I_2$.
\end{problem} 
  
\begin{solution}
This can be proved through a reduction. Firstly, note the construction of $\I$:

\begin{equation*}
    S_1(k_1, m) \oplus S_2(k_2, m) = \I_1 \oplus \I_2
\end{equation*}

\noindent By logic and the laws of XOR, we are essentially just proving $(\I_1 \rightarrow \I) \lor (\I_2 \rightarrow \I)$. So by proving that $\I_1$ is secure, we prove $\I$ is secure.
The proof for $\I_2$ will be the same thing.

\subsection*{Game}
Let $\A$ be the efficent adversary attacking $\I$, assumed to have some adv $\epsilon$. Then, let $\B$ be another efficent adversary, attacking $\I_1$.
Assume the contrapositive: If $\I$ is insecure, then $\I_1$ is insecure.

\begin{itemize}
    \item $\B$ can query $S_1(k_1, \cdot)$. $\A$'s query will be intercepted by $\B$.
    \item Before anything, $\B$ will choose rnd $k_2 \to \K_2$.
    \begin{itemize}
        \item $\B$ is now essentially simulating $\I_2$. This is possible, as we assume that we know how $\I_2$ works.
    \end{itemize}
    \item Whenever $\A$ makes a query with $m$:
    \begin{itemize}
        \item $\B$ forwards query to Challenger, gets $t_1 = S_1(k_1, m)$.
        \item $\B$ then computes $t_2 = S_2(k_2, m)$.
        \item $\B$ computes $t = t_1 \oplus t_2$.
        \item $\B$ sends $t$ to $\A$
    \end{itemize}
    \item $\A$ will output some forgery $(m^*, t^*)$. $\B$ intercepts this.
    \item $\B$ then computes:
    \begin{itemize}
        \item $t^*_2 = S_2(k_2, m^*)$
        \item $t^*_1 = t^* \oplus t^*_2$
    \end{itemize}
    \item $\B$ returns forgery $(m^*, t^*_1)$.
\end{itemize}

\subsection*{Analysis}
\paragraph{Validity} If $\A$ wins, by initial assumption, then $(m^*, t^*)$ is a valid forged tag. By the definition of $\I$:
\begin{equation*}
    t^* = S_1(k_1, m^*) \oplus S_2(k_2, m^*)
\end{equation*}

\noindent Therefore what $\B$ computes is then:
\begin{equation*}
    \begin{aligned}
        t^*_1 &= t^* \oplus S_2(k_2, m^*)\\
        &= (S_1(k_1, m^*) \oplus S_2(k_2, m^*)) \oplus S_2(k_2, m^*)\\
        &= S_1(k_1, m^*)
    \end{aligned}
\end{equation*}

\noindent By the initial assumption, this means that $\B$'s forgery will also then be valid.

\paragraph{Freshness} If $\A$ wins, $m^*$ is never sent to $\I$. $\B$ always forwards $\A$'s queries to $\I_1$. Therefore, if $\A$ has won, then $\B$ hasn't
sent any new message $m^*$ to $\I_1$. Therefore $\B$'s forgery is a valid existential tag for $\I_1$.

\subsection*{End}
Since, when $\A$ wins, $\B$ also wins, the advantage of $\B$ will be at least that of $\A$:
\begin{equation*}
    \MACadv[\B, \I_1] \geq \MACadv[\A, \I] 
\end{equation*}

\noindent By our initial assumption, then $\B$ will therefore then be non-negligible. This proves the contrapositive, therefore proving the OG statement. $\qed$
\end{solution}



\ifsol\else\newpage\fi

\begin{problem} (45 points)
For this problem, it may be helpful to review discussion 7 and section 9.4.2 in the textbook.

    Consider the following candidate authenticated encryption scheme. The idea is to apply a MAC to a plaintext, add a checksum value to detect decryption errors\footnote{Of course, verifying the MAC also suffices to detect decryption errors. Adding the checksum is an attempt at redundancy.}, and then encrypt the plaintext, MAC, and checksum value together.

    Let $\E=(E,D)$ be a block cipher with  $\M = \K = \C = \left\{0,1 \right\}^{128}$, and $\I=(S,V)$ be a MAC with tag space $\left\{0,1 \right\}^{128}$. For the checksum, we'll use a function $\checksum : \left\{0,1 \right\}^* \to \left\{0,1 \right\}^{8}$ that outputs a one-byte tag. 

    Define the scheme $\E'  = (E',D')$ as follows:

    \begin{figure}[h!]
        \center
        \begin{subfigure}{0.45\linewidth}
            \titlecodebox{$E'((k_e, k_m), m)$: }{
                $t \gets S(k_m, m)$\\
                $x \gets \checksum(m \| t) $\\
                $c \gets $  encryption of $m \| t \| x \| 0^{120} $ under $k_e$, \\
                \> using $E$ in randomized CBC mode\\
                Return $c$
            }
        \end{subfigure}
        \begin{subfigure}{0.45\linewidth}
            \titlecodebox{$D'((k_e,k_m), c)$}{
                $(m \| t \| x\| 0^{120}) \gets $ decryption of $c$ under $k_e$, \\
                \> using $D$ in randomized CBC mode\\
                If $x \neq \checksum(m \| t)$:\\
                \> Output \textsf{reject}\\
                Elif $V(k_m, m, t) = \textsf{reject}:$ \\
                \> Output \textsf{reject}\\
                Else:\\
                \> Output $m$
            }
        \end{subfigure}
    \end{figure}

Notice that $E'$ pads the plaintext with 120 bits of zeros. You may assume that all message sizes are multiples of 128 bits. This ensures that the second-to-last block of every ciphertext is an encryption of the tag, and the last block is an encryption of the checksum followed by 120 zeros.

Prove that $\E'$ is not AE-secure.
\end{problem}



\begin{solution}
This protocol's main issue is with WHERE the checksum check is placed for the decryption. At decryption, checksum is verified \textbf{before} the MAC, allowing an adversary to learn information about the plaintext before the MAC.

\subsection*{Idea}
Note two key properties regarding this question:
\begin{itemize}
    \item \textbf{Checksum Before MAC:} $D'$ checks $x = \checksum(m \| t)$ before verifying $V(k_m, m, t)$. This means that a REJECT at the checksum stage reveals different information than a REJECT at the MAC stage.
    \item \textbf{CBC Malleability:} Then, by the malleability of CBC blocks, it is possible to alter the ciphertext block $C_{i-1}$ such that it XORs the following modification into plaintext block $P_i$:
    \begin{equation*}
        P_i' = D(k_e, C_i) \oplus (C_{i-1} \oplus \Delta) = P_i \oplus \Delta
    \end{equation*}
\end{itemize}

\subsection*{Game}
Let $\A$ be an efficent adversary atking $\E'$:
\begin{itemize}
    \item When $\A$ requests an encryption of message $m$, it gets $c = (C_0, C_1, \ldots, C_{n-1}, C_n)$. It will then do the following:
    \begin{itemize}
        \item $C_{n-1}$ encrypts the MAC tag $t$.
        \item $C_n$ encrypts the checksum $x \| 0^{120}$.
    \end{itemize}
    \item $\A$ will then make $c' = (C_0, C_1, \ldots, C_{n-1}, C_{n-1})$ (just replace the last block with the second-to-last block).
    \item Now, $\A$ will send $c'$ to the decryption oracle.
    \item Now, what $\A$ will do will depend on the timing of the decryption function:
    \begin{itemize}
        \item If the code sends REJECT at \textbf{checksum check} stage, $\A$ will output some guess about plaintext.
        \item If a REJECT is sent at the \textbf{MAC check} stage (or accepts), $\A$ will then output a different guess.
    \end{itemize}
    \item $\A$ should run this entire game (with same message $m$) multiple times, in order to get fresh encs, with different IVs, to increase confidence.
\end{itemize}

\subsection*{Analysis}
\paragraph{$c'$} When $c'$ is decrypted, the last plaintext block becomes:
\begin{equation*}
    P_n' = D(k_e, C_{n-1}) \oplus C_{n-2} = t \oplus C_{n-2}
\end{equation*}

\noindent CHECKSUM only passes iff the first byte of $P_n' == \checksum(m \| t)$. By the definition of $\E'$, this should only occur w/ probability $1/256$ (per query).

\paragraph{Leak} If the oracle accepts (or reaches MAC verification), $\A$ then learns that:
\begin{equation*}
    \text{First byte of } (t \oplus C_{n-2}) = \checksum(m \| t)
\end{equation*}

\noindent Since $t = S(k_m, m)$ depends on $m$, this leaks info about the plaintext. By def of $\E'$, by around 256 queries, $\A$ should be able to recover one byte of $m$ with high probability:
\begin{equation*}
    \Pr[\text{success}] = 1 - \left(\frac{255}{256}\right)^{256} \approx 0.63
\end{equation*}

\paragraph{CCA-insecure} By repeating this process for different byte positions (9.4.2 describes this as just doing msg shifting iteratively), $\A$ should then be able to recover the entire plaintext $m$ eventually. This breaks IND-CCA security, implying that AE-security is also broken.

\subsection*{End}
Since $\A$ can recover plaintext info with non-negligible adv using polynomial queries:
\begin{equation*}
    \textsf{AEadv}[\A, \E'] \geq \frac{1}{256}
\end{equation*}

\noindent This probability is \textbf{not negligible} (it is a constant). Therefore, $\E'$ is not AE-secure. $\qed$
\end{solution}



\end{document}
